\chapter{Track fitting}
\label{ch:fitting}

Track fitting turns the clustered track candidates of the pattern recognition into kinematics: a kinetic energy, a
polar and an azimuthal angle, and a vertex for every fitted track. This chapter describes the FairRoot task that drives
the fit (\code{AtFitterTask}), the fitter interface it owns (\code{EventFit::AtFitter}), the two maintained fitters
(the GENFIT Kalman fitter \code{EventFit::AtGenfitter} and the Unscented Kalman Filter \code{EventFit::AtFitterUKF}),
the remaining fitters in \url{AtReconstruction/AtFitter/}, and the output classes. It closes with the pitfalls that
have been measured on a1975 data. It supersedes the stand-alone manual
\url{docs/UKF_GENFIT_manual/ukf_genfit_manual.tex}.

Installation of GENFIT and the environment needed to run the macros are covered in Chapter~\ref{ch:installation}; the
classes \code{AtTrack}, \code{AtHitCluster} and \code{AtPatternEvent} in Chapter~\ref{ch:datamodel}; the PID
estimators (\code{AtSpyralPID}, \code{AtHelix3D}) and the drawing of gates in Chapter~\ref{ch:pid}.

\section{Overview}
\label{sec:fit-overview}

The fitter reads only the branch \code{AtPatternEvent}. Everything upstream (unpacking, PSA, space-charge correction,
pattern recognition) is therefore run once and written to a ``reco'' file, and the fit is repeated on that file as
often as needed with different settings. All a1975 production scripts work this way: the input of the fit is a file
\code{<run>\_reco.root} (or \code{<run>\_multifit\_reco.root}), and the output is a new file holding
\code{AtTrackingEvent} and, when an \code{AtPIDTask} runs in the same job, \code{AtPIDEvent}.

Table~\ref{tab:fit-fitters} lists the fitting classes present on the branch.

\begin{table}[h]
\centering
\small
\begin{tabular}{@{}p{4.6cm}p{3.0cm}p{7.2cm}@{}}
\toprule
Class & Driven by & Status and use \\
\midrule
\code{EventFit::AtGenfitter} & \code{AtFitterTask} & GENFIT \code{KalmanFitterRefTrack} with \code{RKTrackRep}. The
GENFIT fitter of the a1975 and a1954 productions. Compiled only when GENFIT2 is found. \\
\code{EventFit::AtFitterUKF} & \code{AtFitterTask} & Unscented Kalman Filter on the vendored OpenKF library, CATIMA
energy loss, RTS smoother. Compiled only when Eigen3 is found. \\
\code{EventFit::\allowbreak AtFitterUKFMulti} & \code{AtFitterTask} & Runs several \code{AtFitterUKF} hypotheses per track and
keeps one. \\
\code{AtFITTER::AtGenfit} & \code{AtFitterTaskOld} & Original GENFIT fitter, marked \code{[[deprecated]]}; kept to
reproduce older results. \\
\code{MCFitter::AtMCFitter}, \code{AtMCFission} & \code{AtMCFitterTask} & Monte Carlo fitters that simulate and
digitise trial events; \code{AtMCFission} is the fission implementation. \\
\code{MCMinimization::\allowbreak AtMCQMinimization} & a macro & Monte Carlo minimisation on the pad charge, ported from the
FairRootv18.00 branch. Not a task. \\
\bottomrule
\end{tabular}
\caption{Fitting classes on the OpenKF-Claude branch.}
\label{tab:fit-fitters}
\end{table}

\code{docs/development/fitting-status.md} states that fitting on this branch is not a settled architecture and that
the code is the reference. Every default quoted below was read from the sources at the time of writing.

\section{The fitter interface and AtFitterTask}
\label{sec:fit-task}

\subsection{EventFit::AtFitter}

\code{EventFit::AtFitter} (\url{AtReconstruction/AtFitter/AtFitter.h}) is a \code{TObject} with two pure virtual
methods, \code{Init()} and the protected \code{GetFittedTrack(AtTrack*, AtFitMetadata*, AtRawEvent*, AtEvent*)}, and a
virtual \code{FitEvent()}. The default \code{FitEvent()} copies the track candidates out of the \code{AtPatternEvent},
stores the copies in the \code{AtTrackingEvent} (\code{SetTrackArray}), calls \code{GetFittedTrack()} once per
candidate, and adds every non-null result with \code{AddFittedTrack()}. A fitter that returns \code{nullptr} for a
track drops it from the output; there is no record of the rejection apart from the missing track ID.

A new fitter derives from \code{AtFitter}, implements \code{GetFittedTrack()} and returns a heap-allocated
\code{AtFittedTrack} that the event takes ownership of. It may override \code{FitEvent()} when it needs to work on the
whole event (\code{AtGenfitter} does so for continuity merging, Section~\ref{sec:gf-merge}).

\subsection{AtFitterTask}

\code{AtFitterTask} (\url{AtReconstruction/AtFitterTask.h}) is the FairRoot task. Its constructor takes ownership of
the fitter:
\begin{verbatim}
AtFitterTask(std::unique_ptr<EventFit::AtFitter> fitter);
\end{verbatim}

\begin{table}[h]
\centering
\small
\begin{tabular}{@{}p{5.6cm}p{2.8cm}p{6.4cm}@{}}
\toprule
Setter & Default & Effect \\
\midrule
\code{SetInputBranch(TString)} & \code{AtPatternEvent} & Branch holding the \code{AtPatternEvent}. \\
\code{SetOutputBranch(TString)} & \code{AtTrackingEvent} & Branch the \code{AtTrackingEvent} is registered as. \\
\code{SetPersistence(Bool\_t)} & \code{kFALSE} & Write the output branches to the file. \\
\code{SetFitMetadataBranch(TString)} & unset & Register an \code{AtFitMetadata} branch (persistent only together
with \code{SetPersistence}). \\
\code{SetRawEventBranch(TString)} & unset & Pass the \code{AtRawEvent} to the fitter. \\
\code{SetEventBranch(TString)} & unset & Pass the \code{AtEvent} to the fitter. \\
\bottomrule
\end{tabular}
\caption{\code{AtFitterTask} configuration.}
\end{table}

Points about the task that matter in practice:
\begin{itemize}
\item \code{SetParContainers()} requests \code{AtDigiPar} and aborts with a fatal log message if it is missing, so
  every fit macro opens a parameter file even when the fitter uses none of its values. The a1975 macros pass the par
  file the reconstruction ran with (\code{parName} argument of \code{fitGenfitter\_a1975\_deuterium.C}); the macro
  header records that a hard-coded par once made a drift-velocity study fit every variant against the baseline.
\item The task never calls \code{fFitter->Init()}. \code{AtGenfitter} initialises itself on the first track (at which
  point the geometry is loaded), and \code{AtFitterUKF} creates its filter lazily on the first track.
\item \code{Exec()} copies the event ID of the \code{AtPatternEvent} into the \code{AtTrackingEvent}
  (\code{SetEventID}). Use it, not the tree index, to join fit results with other data: a fit made on a pre-gated
  file renumbers its entries from zero.
\item If the input \code{TClonesArray} is empty, \code{Exec()} returns before clearing its output arrays. The tree
  entry written for that event is the \emph{previous} event's \code{AtTrackingEvent}; see
  Section~\ref{sec:fit-pitfalls}.
\end{itemize}

\section{Frames, vertex end and field sign}
\label{sec:fit-frames}

\paragraph{Digi and lab frames.} Hit and cluster positions in \code{AtPatternEvent} are in the digitisation frame, where
$z_\mathrm{digi}$ comes from the drift time. Both maintained fitters convert to a lab frame with the pad plane at a
configurable position:
\[
  z_\mathrm{lab} = z_\mathrm{pad} - z_\mathrm{digi},\qquad x_\mathrm{lab}=x_\mathrm{digi},\qquad
  y_\mathrm{lab}=y_\mathrm{digi},
\]
with $z_\mathrm{pad}$ set by \code{SetZPadPlane()} (default 1000\,mm in \code{AtGenfitter}; in \code{AtFitterUKF} the
default is $-1$, which disables the conversion, and all macros set 1000). The fitted angles, vertex and smoothed
positions are returned in this lab frame. The vertex end of a beam-induced track is the end with the highest
$z_\mathrm{digi}$, which is the lowest $z_\mathrm{lab}$.

\paragraph{Simulated tracks.} A simulated track is reconstructed in the same digitisation frame as the data, so the
conversion above applies to it unchanged. With a fit sign that matches the generated field, the fitted polar angle
compares directly with the Monte Carlo angle (\url{macro/Simulation/ATTPC/16C_dt/sim_fit_sign_check.C} prints both
$\theta-\theta_\mathrm{true}$ and $180^\circ-\theta-\theta_\mathrm{true}$, so that the relation can be measured for a new
simulation). An earlier a1975 simulation gave a different result, described in Chapter~\ref{ch:conventions}.

\paragraph{Field sign.} \code{AtDigiPar} carries only the magnitude of the field. Both fitters take a signed field, and
the sign selects the helix handedness the fit can follow. With the wrong sign a fit does not fail loudly; it converges
to a wrong momentum. The measured cases in the repository are:

\begin{table}[h]
\centering
\small
\begin{tabular}{@{}p{3.6cm}p{3.2cm}p{8.0cm}@{}}
\toprule
Case & Field passed & Source of the statement \\
\midrule
a1975 data, UKF & $-2.85$\,T & \url{macro/Unpack_HDF5/a1975/UKF/README.md}: on run\_0116 the median $\chi^2/\mathrm{ndf}$
is 0.04 and the median KE 7.9\,MeV at $-2.85$\,T, against 3.09 and 152\,MeV at $+2.85$\,T. \\
a1975 data, \code{AtGenfitter} & $-2.85$\,T & Default of the constructor and of every a1975 fit macro and production
script. \\
16C(p,p) simulation, UKF & $+2.85$\,T & \url{macro/Simulation/ATTPC/16C_pp/run_ukf_only.C}; the Geant4 field of
\code{C16\_pp\_sim.C} is $+28.5$\,kG. \\
46Ar(3He,d) simulation, \code{AtGenfitter} & $-2.85$\,T & Header of
\url{macro/Simulation/ATTPC/46Ar_3Hed/fitGenfitter_Ar46.C}: at $+2.85$\,T backward tracks at
$\theta_\mathrm{lab}=110$--$125^\circ$ collapse to 1\,\% of the true energy and 94.1\,\% of truth tracks are fitted,
against 98.2\,\% at $-2.85$\,T. That simulation is generated with a Geant4 field of $-28.5$\,kG. \\
a1975 D$_2$ simulations, \code{AtGenfitter} & $-2.85$\,T & \url{macro/Simulation/ATTPC/16C_dt/sim_fit_sign_check.C}, 1500
events at gain 100\,000: 97.7\,\% of tritons fitted at $-2.85$\,T against 88.3\,\% at $+2.85$\,T; proton
$\ln(\mathrm{KE}/\mathrm{KE}_\mathrm{true})$ at the 16th and 84th percentiles $-0.015$ and $+0.010$ against $-0.089$ and
$+0.040$. Generated with a Geant4 field of $-28.5$\,kG. \\
\code{AtMCQMinimization}, a1975 run\_0106 & \code{SetBFieldSign(-1)} &
\url{macro/tests/AT-TPC/MCQMinimizationOnData.C}. \\
\bottomrule
\end{tabular}
\caption{Field signs that have been measured.}
\label{tab:fit-bsign}
\end{table}

The simulations of the a1975 (d,t) channel (\code{C16\_dt\_sim.C}) are generated with $-28.5$\,kG ``in the data
convention'', according to the comment in that macro. The sign is therefore a property of the pair (simulation field,
fitter), and must be checked on a new setup before anything else: fit a few hundred events with both signs and compare
the elastic locus or, on simulation, the ratio to truth as a function of $\theta_\mathrm{lab}$. Forward tracks can look
identical with both signs (the 46Ar table above is identical below $90^\circ$); the backward tracks decide.

\section{AtGenfitter}
\label{sec:gf}

\code{EventFit::AtGenfitter} (\url{AtReconstruction/AtFitter/AtGenfitter.h}, \code{.cxx}) was written as a replacement
for the legacy \code{AtFITTER::AtGenfit}. It keeps the GENFIT core (\code{genfit::KalmanFitterRefTrack},
\code{genfit::RKTrackRep}, measurements built by \code{genfit::AtSpacepointMeasurement} from \code{AtHitCluster}) and
drops the merging, reclustering and single-vertex logic of the old class. It derives from the modern
\code{EventFit::AtFitter}, so it runs in \code{AtFitterTask}.

\subsection{Construction and particle hypothesis}

\begingroup\footnotesize
\begin{verbatim}
AtGenfitter(Double_t bFieldTesla = -2.85, Int_t pdg = 1000010020,
            Double_t massAmu = 2.0135532, Int_t Z = 1,
            std::string eLossFile = "",
            Bool_t noMatEffects = kTRUE, Int_t minIter = 2, Int_t maxIter = 5);
\end{verbatim}
\endgroup

\begin{itemize}
\item \code{bFieldTesla} is the signed field along $z$ in tesla. \code{Init()} hands GENFIT a
  \code{genfit::ConstField(0, 0, 10*bFieldTesla)} in kilogauss.
\item \code{pdg}, \code{massAmu} and \code{Z} define the hypothesis. The PDG code goes to \code{RKTrackRep} and to the
  measurement candidate; \code{Init()} adds deuteron (1000010020) and triton (1000010030) entries to
  \code{TDatabasePDG} when they are missing. The mass (in u) is used to convert the fitted momentum into kinetic
  energy and for the range constraint. \code{SetParticle(pdg, massAmu, Z)} overrides all three after construction.
  The species is fixed per job: to fit a second hypothesis, run a second job.
\item \code{eLossFile} is an ASCII \dedx{} table; empty means no table (Section~\ref{sec:gf-eloss}).
\item \code{noMatEffects = kTRUE} (the default) switches GENFIT's material effects off: no energy loss, no multiple
  scattering, no straggling. The a1975 macros take a \code{matEffects} argument and pass \code{!matEffects}.
\item \code{minIter}, \code{maxIter} are the iteration limits of \code{KalmanFitterRefTrack}
  (\code{SetIterations(mn, mx)} changes them).
\end{itemize}

GENFIT needs the geometry: \code{Init()} builds a \code{genfit::TGeoMaterialInterface}, which navigates
\code{gGeoManager}. The fit macros load it explicitly before the run starts:
\begin{verbatim}
   // genfit material/propagation needs gGeoManager loaded before the fit
   if (gROOT->FindObject("FAIRGeom") == nullptr) {
      TFile *gf = TFile::Open(geoManFile);
      gf->Get("FAIRGeom");
   }
\end{verbatim}
With material effects on, the medium of that geometry \emph{is} the material model. The comment in
\code{fitGenfitter\_a1975\_deuterium.C} records that the default geometry \code{ATTPC\_H1bar\_\allowbreak geomanager.root} contains
ordinary hydrogen; for the D2 runs the productions pass \code{ATTPC\_D300torr\_v2\_\allowbreak geomanager.root} (medium
\code{TargetD2\_300}, $A=2.0141$, $Z=1$, $6.564\times10^{-5}$\,g/cm$^3$).

\code{MaterialEffects} and \code{FieldManager} are GENFIT singletons. \code{Init()} configures them for the instance
being initialised, so two \code{AtGenfitter} instances with different field or material settings in one process
share whatever the last initialised one set. Run different configurations in separate jobs.

\subsection{What happens to one track}
\label{sec:gf-steps}

\code{GetFittedTrack()} performs the following steps, in this order.
\begin{enumerate}
\item Reject the candidate if it has fewer than \code{SetMinClusters(n)} clusters (default 4).
\item If any option needs it (PID gate, Spyral seed, Spyral or 3D direction), compute the \code{AtSpyralPID} estimate
  and, for the 3D options, the \code{AtHelix3D} estimate. If a PID gate is set and the track lies outside it, return
  \code{nullptr} (Section~\ref{sec:gf-pid}).
\item Convert cluster positions to the lab frame and, unless \code{SetUseClusterOrder(kTRUE)} is set
  (Section~\ref{sec:gf-order}), sort them by $z_\mathrm{lab}$.
\item Reject the track if no cluster lies within \code{SetVertexAxisMaxDist} (default 150\,mm) of the beam axis.
\item Decide the vertex end. By default it is the lowest-$z_\mathrm{lab}$ cluster, with the seed direction taken from
  that cluster to the next one. With \code{SetBackwardSeedFix(kTRUE)} a track judged backward is seeded from the
  highest-$z_\mathrm{lab}$ end instead (Section~\ref{sec:gf-direction}).
\item Seed the momentum: $B\rho = |B|\,r/|\sin\theta|$ with $r$ the PRA circle radius (\code{GetGeoRadius()}) and
  $\theta$ the polar angle of the two-cluster direction, $p = 0.299792458\,B\rho$ GeV/$c$ for unit charge. Tracks
  with a seed above 100\,GeV/$c$ or non-positive are rejected. \code{SetSeedFromSpyral(kTRUE)} replaces $B\rho$ and
  $\theta$ by the \code{AtSpyralPID} values. The seed covariance is 1\,mm$^2$ in position and $(0.3p)^2$ in each
  momentum component.
\item Order the measurements along the direction of travel (reversed for a backward seed), apply truncation if set,
  and fill the \code{genfit::TrackCand} with one \code{AtSpacepointMeasurement} per cluster.
\item Optionally compute a range constraint (Section~\ref{sec:gf-range}).
\item Re-seed the shared TGeo navigator (Section~\ref{sec:gf-navigator}) and fit with
  \code{processTrackWithRep()}. Read the state at the first measurement point; if
  \code{SetBackExtrapToAxis(kTRUE)}, extrapolate it to the beam axis.
\item If FIND is on and the full-length fit fails its threshold, refit prefixes (Section~\ref{sec:gf-find}).
\item If the fit threw, material effects were on and the fallback is enabled, refit without material effects.
\item Add the manual vertex-gap energy loss if \code{SetManualELoss} was called.
\item Drop the track if the fitted polar angle lies outside \code{SetThetaWindow(min, max)} (default 10--170$^\circ$).
\item Fill the \code{AtFittedTrack} and its \code{AtFitTrackMetadata} (Section~\ref{sec:fit-output}).
\end{enumerate}
A track that throws inside GENFIT and is not rescued by the fallback is dropped. A fit that returns but with
$\mathrm{ndf}\le 0$ is \emph{kept}; its kinematics are not meaningful (Section~\ref{sec:fit-pitfalls}).

\subsection{Measurement model and the meaning of $\chi^2$}
\label{sec:gf-meas}

Each cluster becomes a 3D space point with a diagonal covariance, set by \code{AtGenfitter} in mm$^2$ and converted to
cm$^2$ by \code{AtSpacepointMeasurement}:
\[
  \sigma_x^2=\sigma_y^2 = s\,(\sigma_0^2 + D_T^2 L),\qquad
  \sigma_z^2 = s\,(f_z\sigma_0^2 + D_L^2 L),
\]
where $\sigma_0$ is \code{SetMeasSigma(mm)} (default 4.0\,mm), $L$ the drift distance in cm
($(z_\mathrm{pad}-z_\mathrm{lab})/10$), $D_T, D_L$ the coefficients of \code{SetDiffusion(transv, long)} in
mm$/\sqrt{\mathrm{cm}}$ (default 0, 0), $f_z$ \code{SetZLongFactor} (default 2.0) and $s$ the optional charge scale
$\mathrm{clamp}(q_\mathrm{ref}/Q, 0.25, 4)$ enabled by \code{SetChargeCovRef(qRef)} with $q_\mathrm{ref}>0$ (default
0, off). With the defaults this reduces to $(\sigma_0^2,\sigma_0^2,2\sigma_0^2)$.

The 4\,mm default is far larger than the real residuals. The comments in \code{AtGenfitter.h} and in
\code{fitGenfitter\_Ar46.C} quote the residuals implied by $\chi^2/\mathrm{ndf}$ as about 0.6\,mm, and the median
$\chi^2/\mathrm{ndf}$ of an a1975 (d,p) production at 4\,mm as 0.09 (1.48 for backward tracks).
\url{macro/Unpack_HDF5/a1975/UKF/pipeline/GENFIT_PROGRESS.md} records the trade on run\_0106 with material effects off:
4\,mm gives a median $\chi^2/\mathrm{ndf}$ of 0.105 and 79\,\% good fits, 1.5\,mm gives 0.75 and 58\,\%, 1.0\,mm gives
1.69 and 54\,\%; the yield is lost because the larger $\sigma_0$ absorbs multiple scattering that the model without
material effects cannot follow. Consequences:
\begin{itemize}
\item $\chi^2/\mathrm{ndf}$ at 4\,mm is not a goodness-of-fit in the statistical sense. Its scale grows with the number
  of turns and the length of the track, so a fixed cut selects on track length. The 46Ar macro header measured that the
  production cut of 5 removed 29\,\% of the longest spirals there, and changed its own default to 0.6\,mm.
\item Thresholds (the quality cut, FIND) have to be set from the measured distribution of the production at hand,
  not carried over from another $\sigma_0$.
\item The $p$-value is not computed: \code{AtGenfitter} stores 0.
\end{itemize}

\subsection{Iterations, quality flag and angular window}

\code{KalmanFitterRefTrack} iterates between \code{minIter} and \code{maxIter}; \code{GetFitConverged()} in the metadata
is GENFIT's \code{isFitConverged()}. Two post-fit criteria are applied:
\begin{itemize}
\item \code{SetThetaWindow(minDeg, maxDeg)}, default 10--170$^\circ$: the track is \emph{dropped} if the polar angle of
  the final momentum (after any extrapolation) lies outside.
\item \code{SetQualityCut(chi2NdfMax, keMin, keMax)}, default (5.0, 0.5, 60): a \emph{flag} only.
  \code{GetGoodFit()} is true when the fit converged, $\mathrm{ndf}>0$, $\chi^2/\mathrm{ndf}<$ \code{chi2NdfMax} and
  \code{keMin} $<$ KE $<$ \code{keMax} (KE after extrapolation). Failing tracks stay in the output.
\end{itemize}

\subsection{Material effects and energy loss}
\label{sec:gf-eloss}

With material effects off, GENFIT transports the state geometrically: the fitted momentum is that of a helix through
the clusters, with no energy loss. This is fast and was the production setting of the early a1975 fits. With
material effects on, the propagation applies energy loss, multiple scattering and straggling from the material of the
geometry. \code{Init()} always switches bremsstrahlung off and calls
\code{ignoreBoundariesBetweenEqualMaterials(false)}; the comment there records a factor 10 in CPU time on a1954
run\_0055 with bit-identical output, because the AT-TPC active volume is one homogeneous gas volume.

\paragraph{GENFIT's own energy loss and the $\beta\gamma<0.05$ hole.} GENFIT's \code{MaterialEffects::\allowbreak dEdx()} uses
Bethe--Bloch only for $\beta\gamma\ge 0.05$. Below that it calls \code{dEdxParam()}, which evaluates a loaded \dedx{}
curve and returns 0 when there is none (\url{GenFit/trackReps/src/MaterialEffects.cc}). Without a table, material
effects therefore apply \emph{no} stopping power below $\beta\gamma = 0.05$: KE $<1.17$\,MeV for a proton and
KE $<3.5$\,MeV for a triton. For (d,t) at a1975, where the triton low branch is 0.8--6\,MeV, this is most of the path
of most tracks. Three ways to fill the hole exist; they are listed from oldest to adopted.

\paragraph{Table in parameterisation-only mode.} Passing \code{eLossFile} to the constructor loads the table with
\code{setEnergyLossFile()} and, unless the hybrid mode is set, calls \code{useEnergyLossParam()}, which makes the
table the only source of \dedx{}. In that mode GENFIT throws a fatal exception for every step above the highest
energy of the table. The comment in \code{Init()} explains that this, when called without a table, truncated the
reference track of about 87\,\% of tracks and was the origin of the belief that material effects were unusable.

\paragraph{Hybrid table.} \code{SetELossHybrid(on, rho)} loads the table without parameterisation-only mode:
the table serves only $\beta\gamma<0.05$ and Bethe--Bloch the rest. \code{dEdxParam()} multiplies the curve
(MeV/(mg/cm$^2$)) by one global density, which is why the setter takes the gas density in g/cm$^3$ (stored internally
in mg/cm$^3$ and passed to \code{setGasMediumDensity()}); the GENFIT default is 0.15328\,mg/cm$^3$. Tables are generated
per species, gas and pressure by \code{make\_eloss\_table.C}; the header of \code{prod\_dp\_clord.sh} gives the
call used for the proton table,
\begin{verbatim}
make_eloss_table.C("proton_D2_300torr.txt", 1, 1, 1.00782503207,
                  6.61e-5, 2.014)
\end{verbatim}
and \code{prod\_dt\_catima.sh} records that the triton table was regenerated with the deuterium mass 2.014 instead of
2, a flat $+0.70$\,\% on every point.

\paragraph{CATIMA in GENFIT.} The GENFIT installation used on this branch has an optional CATIMA backend, compiled
with \code{-DGENFIT\_USE\_CATIMA=ON} (option in GENFIT's \code{CMakeLists.txt}). Without it the setters below compile
and do nothing. \url{macro/Unpack_HDF5/C15d/README.md} gives the check
\code{ldd \$GENFIT/lib/libgenfit2.so | grep catima}.
\begin{itemize}
\item \code{SetCatimaMaterial(msc, straggling)} replaces GENFIT's Highland multiple scattering and Bohr/Urban
  straggling by CATIMA's angular and energy variances per unit thickness. Default (off, off).
\item \code{SetCatimaELoss(on, fullRange = kFALSE)} takes \dedx{} from CATIMA with the material of the step. With
  \code{fullRange = kFALSE} only the $\beta\gamma<0.05$ branch is replaced; \code{kTRUE} replaces Bethe--Bloch too.
  If CATIMA cannot supply a value, the table (if loaded) is the per-step fallback, so the productions keep the table
  loaded as well.
\end{itemize}
The production script \url{macro/Unpack_HDF5/a1975/D2_UKF/prod_dt_catima.sh} documents the measurements behind the
adopted setting. GENFIT's own material models left 62\,\% of triton fits with $\mathrm{ndf}<0$; with CATIMA MSC,
CATIMA straggling and the table, 3\,\%; removing any one of the three brings the collapse back to about 50\,\%.
Replacing the table by CATIMA \dedx{} below $\beta\gamma=0.05$ reduced collapsed fits further from 2.53\,\% to
0.94\,\% over 8 runs and did not move the 3.103\,MeV peak; the full-range variant changed nothing more.

\paragraph{Fallback without material effects.} \code{SetMatEffectsFallback(on)}, default \code{kTRUE}: when a fit with
material effects throws, the track is refitted with \code{setNoEffects(true)} and kept, flagged
\code{GetMatEffects()==false}, \code{GetMatEffectsFallback()==true}. The two populations come from different models;
mixed in one spectrum they widen it. The a1975 CATIMA productions run with the fallback off, so a failed fit drops out.

\subsection{The shared TGeo navigator}
\label{sec:gf-navigator}

\code{TGeoMaterialInterface} navigates the single \code{gGeoManager} and does not restore its state. A track whose
Runge--Kutta extrapolation wanders leaves the navigator in a state that every later track inherits; the failure is
then contiguous in track order (a job fits its first $N$ tracks and fails almost all after). The comment in
\code{GetFittedTrack()} quotes run\_0016, 5000 events, material effects on, $\sigma_0=4$\,mm: 76.8\,\% of fits
collapsed ($\mathrm{ndf}\le0$) without a fix and 1.8\,\% with it. The fix, now unconditional, re-seeds the navigator
before every track:
\begin{verbatim}
   if (gGeoManager != nullptr) {
      gGeoManager->SetCurrentPoint(0., 0., 50.);
      gGeoManager->SetCurrentDirection(0., 0., 1.);
      gGeoManager->FindNode();
   }
\end{verbatim}
Setting the environment variable \code{ATGENFIT\_TRACE=1} prints one line per track with the navigator node and
material; the probe is read-only. Productions made before this fix show the collapse and should not be compared with
later ones as if only the physics settings differed.

\subsection{The fitted state, back-extrapolation and the vertex gap}
\label{sec:gf-xtr}

\code{genfit::Track::getFittedState()} without arguments returns the state at the \emph{first measurement point}, not
at the reaction vertex. Between the two lies gas in which the ejectile has already lost energy. Two options address it.

\paragraph{\code{SetBackExtrapToAxis(on)}}, default off. Extrapolates the first-point state to the beam axis with
\code{RKTrackRep::extrapolateToLine()} (line through the origin along $z$). If the extrapolation throws, the
first-point state is kept. With material effects on, the extrapolation integrates the energy loss over the gap; with
them off it moves the vertex and leaves $|p|$ unchanged. The setter comment quotes, for a1975 16C(d,t) with the correct
D2 density, the reconstructed g.s.-to-3.103 spacing opening from 2.919 to 3.192\,MeV and the width falling from 0.252
to 0.155\,MeV.

\paragraph{\code{SetManualELoss(density, matA)}.} Adds the energy lost over the vertex gap by hand, from an
\code{AtELossCATIMA} model of the ejectile in hydrogen of mass number \code{matA} (2 for D2; 1 computes the loss in
ordinary hydrogen, twice the stopping at the same mass density). The loss is evaluated at the first-cluster energy and
once more at the corrected energy. It is meant for material effects \emph{off} together with
\code{SetBackExtrapToAxis(kTRUE)}, which supplies the gap length.

\emph{The code does not enforce that.} With material effects on, \code{extrapolateToLine} has already restored the
gap loss and \code{SetManualELoss} adds it a second time. \code{prod\_dt\_catima.sh} measured the cost on the (d,t) low
branch: $+5.25$\,\% on the corrected KE and a 14\,\% wider distribution. In \code{fitGenfitter\_a1975\_deuterium.C} the
argument \code{manualElossDensity} must be 0 whenever \code{matEffects} is \code{kTRUE}; passing 0 disables only the
manual loss (the hybrid table falls back to $6.61\times10^{-5}$\,g/cm$^3$ internally).

\paragraph{The two kinematics slots.} \code{AtGenfitter} fills both:
\begin{itemize}
\item \code{GetKinematics()}: KE, $\theta$, $\phi$ of the fit at the first measurement point, before extrapolation and
  before the manual loss (the ``raw'' fit).
\item \code{GetKinematicsXtr()}: after back-extrapolation and manual loss. With neither option on, the two are equal.
\end{itemize}
The vertex (\code{GetVertex()}) is the position of the final state, i.e.\ the axis point if extrapolated, otherwise the
first fitted point, in lab-frame millimetres. The a1975 D2 analysis macros (for example \code{ex\_dt\_a1975.C}) read \code{GetKinematicsXtr()}. The UKF uses
the slots the other way round (Section~\ref{sec:fit-output}).

\subsection{Track direction and the backward seed}
\label{sec:gf-direction}

The default seed always starts at the lowest-$z_\mathrm{lab}$ end and points the momentum towards $+z_\mathrm{lab}$.
For backward tracks ($\theta_\mathrm{lab}>90^\circ$) that is the stopping end, and GENFIT reflects the track into the
forward hemisphere with a wrong $\theta$ and KE. \code{SetBackwardSeedFix(kTRUE)} (default off) seeds tracks judged
backward from the highest-$z_\mathrm{lab}$ end with the momentum pointing into the track, and reverses the measurement
order so that measurement 0 is again the vertex end. Forward tracks are unchanged. The comment in the code quotes, for
backward protons of 14C(d,p) simulation at 2.85\,T, the median KE ratio to truth going from 0.834 to 0.999 at
92--110$^\circ$ and from 0.893 to 0.999 at 110--140$^\circ$. The (d,p) macro header records the backward yield on
run\_0016 rising from 5.6\,\% to 17.5\,\% of candidates.

Which tracks count as backward is chosen by:
\begin{itemize}
\item default: the PRA polar angle, \code{GetGeoTheta()} $>\pi/2$ (the sign of one RANSAC line in arc length versus
  $z$);
\item \code{SetDirectionFromSpyral(kTRUE)}: the \code{AtSpyralPID} estimate, $\theta_\mathrm{lab}=\pi-$ polar; falls
  back to \code{GetGeoTheta()} when the estimate is invalid. The setter comment records that on a1975 H2 data the two
  disagree on 8\,\% of tracks and the PRA call is the wrong one where the elastic locus can decide;
\item \code{SetDirectionFrom3D(kTRUE)}: the \code{AtHelix3D} majority vote of three indicators (3D rotation sense,
  Spyral direction, Bragg rise) on hits ordered in 3D; on an ambiguous vote the previous source is kept. It takes
  precedence over the Spyral option. The comment quotes 96--97\,\% correct on 16C+d data against 93--95\,\% for
  Spyral, and 100\,\% on simulation truth.
\end{itemize}
All three act only when \code{SetBackwardSeedFix(kTRUE)} is set. The estimators themselves are described in
Chapter~\ref{ch:pid}.

\paragraph{\code{SetSeedFromSpyral(kTRUE)}} (default off) replaces the seed $B\rho$ and polar angle by the
\code{AtSpyralPID} values: a least-squares circle on the first arc and a regression of $\rho$ against $z$ over at
least ten points, instead of the PRA radius and the direction between two adjacent clusters, which on a short tightly
curved track are millimetres apart. The hemisphere fixed by the vertex-end decision is kept.

\subsection{The PID gate}
\label{sec:gf-pid}

\code{SetPIDGate(jsonFile)} loads an \code{AtParticleID} gate (\code{AtParticleID::LoadJSON}) and skips every track
whose PID observables fall outside it, before the fit. The fitter owns its own \code{AtSpyralPID} (built with
$|B|$ of the fitter), so it needs no \code{AtPIDEvent} branch. If the file cannot be loaded, \code{Init()} logs an error
and disables gating; the a1975 macros additionally check that the file exists and print a warning. Either way a mistyped
path fits every track instead of failing.

Options of the gate:
\begin{itemize}
\item \code{SetPIDGateFrom3D(kTRUE)}: test the gate on \code{AtHelix3D}'s $(\sqrt{\mathrm{d}E/\mathrm{d}x}, B\rho)$
  instead of Spyral's. The gate must have been drawn on that plane; the two \dedx{} scales differ.
\item \code{SetDedxGain(gSpyral, g3D)}: multiply the internal \sqdedx{} by per-run gain factors
  (\code{AtTools::AtPIDGainTable}). The gate must have been drawn with the same correction.
  \code{fitGenfitter\_a1975\_deuterium.C} takes a \code{gainTable} argument and aborts if the run is not in it.
\end{itemize}

A gated-out track is absent from \code{AtTrackingEvent}; a later, tighter gate can be applied from \code{AtPIDEvent}
(the PID does not depend on the fit), but a looser one cannot. The (d,p) production header therefore fits with a loose
gate and re-cuts at the analysis stage. The internal \code{AtSpyralPID} runs with the class default minimum number of
points (30); there is no setter on \code{AtGenfitter} to change it, so a gate drawn on observables computed with a
different \code{SetMinPoints} is applied to different numbers inside the fitter (header of
\code{fitGenfitter\_Ar46.C}).

\subsection{Cluster order}
\label{sec:gf-order}

By default the measurements are sorted by $z_\mathrm{lab}$. Drift $z$ is monotonic along a normal spiral, so curling
tracks are sequenced correctly. The assumption fails for tracks close to $90^\circ$, where the advance in $z$ per
cluster is smaller than the $z$ noise and than one time bucket; the code comment gives 0.22\,mm per cluster at
88--92$^\circ$ and 7\,T, and the (d,p) macro 1.42\,mm per cluster against 1.84\,mm per time bucket on a1975 backward
tracks. The sort then orders noise and the filter integrates energy loss along a path the particle never flew.

\code{SetUseClusterOrder(kTRUE)} keeps the order of the cluster array as delivered by the pattern recognition (to be
paired with an ordering clusteriser such as \code{AtPRA::SetUseArcWalk(true)}, Chapter~\ref{ch:patternrecognition}).
It is \emph{not} adopted. The comments in \code{fitGenfitter\_a1975\_deuterium.C} and \code{prod\_dp\_clord.sh} record
a five-track test in which one track lost its fit and one swung by $100^\circ$ in $\theta$, and state that the loss rate
has to be measured on full runs before use.

\subsection{Truncation and the FIND ladder}
\label{sec:gf-find}

\paragraph{Truncation.} \code{SetTruncatePercent(pct)} fits only the first \code{pct}\,\% of the measurement sequence,
counted from the vertex end (0 or $\ge100$ disables it, the default). \code{SetTruncateMinClusters(n)} (default 8) is a
floor. The motivation given in the header: a single circle through a decelerating spiral averages a radius that shrinks
along the track (0.964 to 0.734 from vertex to stopping end on 14C(d,p) at 2.85\,T), biasing the momentum low by about
20\,\% in KE. Only measurements are truncated; the range constraint still uses the full path. $\chi^2$ falls when points
are removed, so a truncation is judged on the excitation-energy spectrum, not on $\chi^2$.

\paragraph{FIND.} \code{SetFindLongest(on, c2Max, stepPct = 5, minPct = 25, maxPct = 95)} refits prefixes of a track
whose full-length fit fails $\chi^2/\mathrm{ndf}<$ \code{c2Max} or throws. The ladder runs from
$\min(100-\mathrm{step},\ \mathrm{maxPct})$ down to \code{minPct} in steps of \code{stepPct}, and stops when a prefix
would have fewer than \code{SetTruncateMinClusters} clusters; for small \code{minPct} that floor is what actually
binds. \code{SetFindMode(mode)} selects:
\begin{itemize}
\item mode 0 (default), longest passing: the first (longest) rung under \code{c2Max} is kept; if none passes, the
  full-length fit is restored. Tracks that pass at full length cost nothing extra.
\item mode 1: if the full-length fit passes it is kept; otherwise every rung is fitted and the longest passing one is
  kept, or, if none passes, the one with the lowest $\chi^2/\mathrm{ndf}$. The winner is refitted at the end.
\end{itemize}
\code{GetLastFindPct()} returns the percentage kept for the last track; it is not written to the output.
\code{c2Max} must come from the measured $\chi^2$ distribution: at $\sigma_0=4$\,mm the median is about 0.09, so a cut
of 5 never fires. Note also that some comments in the header and macro still describe the ladder as ``90/75/50/25'';
the code uses the step ladder above.

\subsection{Range constraint}
\label{sec:gf-range}

\code{SetRangeConstraint(on, gasDensity, matA = 2)} (default off) injects, for tracks that stop, a
\code{genfit::FullMeasurement} at the vertex end whose covariance is tight only along the momentum direction, so that
it constrains $|p|$ from the path length and nothing else. It is a measurement rather than a tight seed because
\code{KalmanFitterRefTrack} inflates the seed covariance between iterations. The range--energy relation comes from an
\code{AtELossCATIMA} model of the ejectile; $\sigma_\mathrm{KE}$ is the range straggling propagated through
$\mathrm{d}R/\mathrm{d}E$ with a fractional floor \code{SetRangeSigmaFloor} (default 0.03).

A track qualifies only if its far end is contained (\code{SetRangeContainment(maxRadius = 240, zMargin = 40,
minLength = 20)} mm) and it passes a Bragg test: the charge per unit path over the last \code{tailFrac} of the path
divided by that over the middle half must exceed \code{minRatio}
(\code{SetBraggTest(minRatio = 1.3, tailFrac = 0.25, minClusters = 10)}; \code{minRatio}~$\le0$ disables the test).
The header explains why: on a1975 run\_0031 with containment alone, 173 of 316 ``contained'' tracks were above 20\,MeV
and were handed about 2.9\,MeV. The counters \code{GetNRangeApplied()}, \code{GetNRangeContained()} and
\code{GetNRangeBraggFail()} report how often each branch fired. \code{prod\_dt\_catima.sh} leaves the range constraint
off.

\subsection{Continuity merging}
\label{sec:gf-merge}

\code{SetMergeContinuity(on, centerDist = 15, radiusFrac = 0.3, gap = 30)} (default off) merges, before fitting, pattern
tracks that are fragments of one trajectory. Two fragments merge when their PRA circle centres are within
\code{centerDist} mm, their radii differ by less than \code{radiusFrac} of the larger, and their closest endpoints (vertex
end = closest to the axis, far end = farthest) are within \code{gap} mm, provided the closest pair is not
vertex-end to vertex-end (two tracks from one vertex). Chains merge transitively. Merging works on a copy of the track
list, so the shared \code{AtPatternEvent} and any other fitter in the job see the original. The merged list is what is
stored as the event's track array. \code{GENFIT\_PROGRESS.md} records the a1975 test: on 30094 (p,d)-gated events the
merge fired 14 times, all neutral.

\subsection{Setter reference}

\begin{table}[p]
\centering
\small
\begin{tabular}{@{}p{6.4cm}p{2.6cm}p{5.8cm}@{}}
\toprule
Setter & Default & Section \\
\midrule
\code{SetParticle(pdg, massAmu, Z)} & ctor & hypothesis \\
\code{SetBField(tesla)} & $-2.85$ & \ref{sec:fit-frames} \\
\code{SetZPadPlane(mm)} & 1000 & \ref{sec:fit-frames} \\
\code{SetMinClusters(n)} & 4 & \ref{sec:gf-steps} \\
\code{SetIterations(min, max)} & 2, 5 & iterations \\
\code{SetVertexAxisMaxDist(mm)} & 150 & \ref{sec:gf-steps} \\
\code{SetMeasSigma(mm)} & 4.0 & \ref{sec:gf-meas} \\
\code{SetDiffusion(DT, DL)} & 0, 0 & \ref{sec:gf-meas} \\
\code{SetZLongFactor(f)} & 2.0 & \ref{sec:gf-meas} \\
\code{SetChargeCovRef(q)} & 0 (off) & \ref{sec:gf-meas} \\
\code{SetThetaWindow(min, max)} & 10, 170 deg & drops tracks \\
\code{SetQualityCut(c2, keMin, keMax)} & 5.0, 0.5, 60 & flag only \\
\code{SetMatEffectsFallback(on)} & \code{kTRUE} & \ref{sec:gf-eloss} \\
\code{SetELossHybrid(on, rho)} & off & \ref{sec:gf-eloss} \\
\code{SetCatimaMaterial(msc, strag)} & off, off & \ref{sec:gf-eloss} \\
\code{SetCatimaELoss(on, full)} & off, off & \ref{sec:gf-eloss} \\
\code{SetBackExtrapToAxis(on)} & off & \ref{sec:gf-xtr} \\
\code{SetManualELoss(rho, matA)} & not set & \ref{sec:gf-xtr} \\
\code{SetBackwardSeedFix(on)} & off & \ref{sec:gf-direction} \\
\code{SetDirectionFromSpyral(on)} & off & \ref{sec:gf-direction} \\
\code{SetDirectionFrom3D(on)} & off & \ref{sec:gf-direction} \\
\code{SetSeedFromSpyral(on)} & off & \ref{sec:gf-direction} \\
\code{SetPIDGate(json)} & none & \ref{sec:gf-pid} \\
\code{SetPIDGateFrom3D(on)} & off & \ref{sec:gf-pid} \\
\code{SetDedxGain(gSpyral, g3D)} & 1, 1 & \ref{sec:gf-pid} \\
\code{SetUseClusterOrder(on)} & off & \ref{sec:gf-order} \\
\code{SetTruncatePercent(pct)} & 0 (off) & \ref{sec:gf-find} \\
\code{SetTruncateMinClusters(n)} & 8 & \ref{sec:gf-find} \\
\code{SetFindLongest(on, c2, step, min, max)} & off; 5, 25, 95 & \ref{sec:gf-find} \\
\code{SetFindMode(mode)} & 0 & \ref{sec:gf-find} \\
\code{SetRangeConstraint(on, rho, matA)} & off; matA 2 & \ref{sec:gf-range} \\
\code{SetRangeContainment(r, zMargin, len)} & 240, 40, 20 & \ref{sec:gf-range} \\
\code{SetRangeSigmaFloor(f)} & 0.03 & \ref{sec:gf-range} \\
\code{SetBraggTest(ratio, tail, nMin)} & 1.3, 0.25, 10 & \ref{sec:gf-range} \\
\code{SetMergeContinuity(on, c, r, gap)} & off; 15, 0.3, 30 & \ref{sec:gf-merge} \\
\bottomrule
\end{tabular}
\caption{\code{AtGenfitter} setters with their defaults (lengths in mm, densities in g/cm$^3$).}
\label{tab:gf-setters}
\end{table}

\section{Worked example: the a1975 16C(d,t) production}
\label{sec:fit-example}

The (d,t) channel of a1975 is fitted by \url{macro/Unpack_HDF5/a1975/D2_UKF/fitGenfitter_a1975_deuterium.C}, driven over
47 runs by \url{macro/Unpack_HDF5/a1975/D2_UKF/prod_dt_catima.sh}. The macro serves every ejectile of the D2 target; its
defaults are the proton of 16C(d,p), and the species is passed as arguments.

\subsection{The macro}

The macro reads \code{<ioDir><run><recoSuffix>.root} and writes \code{<outDir><run>\_genfitter\_\allowbreak <tag><suffix>.root}.
After opening the par file and the geometry (see above), it builds the fitter from its arguments. Condensed, the
statements that matter for the (d,t) production are:
\begingroup\footnotesize
\begin{verbatim}
   TString elossPath = eLossTable.Length() ? (eLossTable.BeginsWith("/") ? eLossTable
                        : dir + "/macro/Unpack_HDF5/a1975/D2_UKF/" + eLossTable) : TString("");
   auto fitter = std::make_unique<EventFit::AtGenfitter>(bField, pdg, massAmu, Z,
                        elossPath.Data(), !matEffects, minIter, maxIter);
   if (catimaMSC || catimaStraggling)
      fitter->SetCatimaMaterial(catimaMSC, catimaStraggling);
   if (eLossTable.Length())
      fitter->SetELossHybrid(kTRUE, manualElossDensity > 0 ? manualElossDensity : 6.61e-5);
   fitter->SetZPadPlane(1000.0);
   fitter->SetMeasSigma(measSigma);
   ...
   fitter->SetMatEffectsFallback(matFallback);
   fitter->SetThetaWindow(thetaMinDeg, thetaMaxDeg);
   fitter->SetBackwardSeedFix(backwardSeedFix);
   fitter->SetBackExtrapToAxis(backExtrap);
   fitter->SetCatimaELoss(catimaELoss, catimaELossFull);
   if (manualElossDensity > 0)
      fitter->SetManualELoss(manualElossDensity, matA);
   ...
   fitter->SetPIDGate(pidGate.Data());       // after checking that the file exists

   AtFitterTask *fitterTask = new AtFitterTask(std::move(fitter));
   fitterTask->SetInputBranch("AtPatternEvent");
   fitterTask->SetOutputBranch("AtTrackingEvent");
   fitterTask->SetFitMetadataBranch("AtFitMetadata");
   fitterTask->SetPersistence(kTRUE);

   AtPIDTask *pidTask = new AtPIDTask();
   pidTask->SetInputBranch("AtPatternEvent");
   pidTask->SetOutputBranch("AtPIDEvent");
   pidTask->SetPersistence(kTRUE);

   run->AddTask(pidTask);
   run->AddTask(fitterTask);
\end{verbatim}
\endgroup
Every option that changes the physics prints a banner to the log (``CATIMA material model: MSC ON, straggling ON'',
``dE/dx from CATIMA \dots''), and the macro warns when \code{catimaELoss} is set without material effects, where it is
inert. The \code{AtPIDTask} in the same job writes the PID of every pattern track, gated or not, so that the analysis
can apply a different gate later.

\subsection{The production script}

\code{prod\_dt\_catima.sh} sources the environment, sets \code{GENFIT} and the library path, and runs one ROOT job per
run through \code{xargs -P}. The call passes every argument positionally:
\begingroup\footnotesize
\begin{verbatim}
  root -l -b -q "fitGenfitter_a1975_deuterium.C(\"${r}_multifit\",-1,\"$REC\",\"\",\"$GF\",\
-2.85,2,5,\"$GATE\",4.0,10.0,170.0,kTRUE,kTRUE,1000010030,3.01550072,1,\"t\",\"_reco\",\
\"ATTPC_D300torr_v2_geomanager.root\",kTRUE,0,2,\"$PAR\",kFALSE,kFALSE,kFALSE,\"$TAB\",\
kTRUE,kTRUE,${CEL},${CFULL})"
\end{verbatim}
\endgroup
Mapped onto the macro arguments, Table~\ref{tab:fit-dt-args} gives the adopted (d,t) configuration.

\begin{table}[h]
\centering
\small
\begin{tabular}{@{}p{3.8cm}p{5.0cm}p{6.0cm}@{}}
\toprule
Argument & Value & Why (from the script header) \\
\midrule
\code{bField} & $-2.85$ & data handedness \\
\code{minIter}, \code{maxIter} & 2, 5 & unchanged from earlier productions \\
\code{pidGate} & \code{pid/triton\_d2\_\allowbreak dv1104.json} & triton gate on the plane of this drift velocity \\
\code{measSigma} & 4.0 & unchanged \\
\code{theta window} & 10--170$^\circ$ & keep backward tracks \\
\code{matEffects} & \code{kTRUE} & the point of the production \\
\code{backwardSeedFix} & \code{kTRUE} & backward seed \\
\code{pdg}, \code{massAmu}, \code{Z}, tag & 1000010030, 3.01550072, 1, \code{t} & triton \\
\code{geoName} & \code{ATTPC\_D300torr\_v2\_\allowbreak geomanager.root} & D2 medium for the material model \\
\code{backExtrap} & \code{kTRUE} & recover the vertex-gap loss \\
\code{manualElossDensity} & 0 & would double count the gap under material effects \\
\code{parName} & \code{ATTPC.a1975\_deuterium\_\allowbreak dv1104.par} & the par of the reconstruction \\
\code{matFallback} & \code{kFALSE} & one physics model in the output \\
\code{eLossTable} & \code{triton\_D2\_300torr.txt} & hybrid table below $\beta\gamma=0.05$ \\
\code{catimaMSC}, \code{catimaStraggling} & \code{kTRUE}, \code{kTRUE} & collapse 62\,\% $\to$ 3\,\% \\
\code{catimaELoss}, \code{catimaELossFull} & \code{kTRUE}, \code{kFALSE} (env.\ \code{CEL}, \code{CFULL}) &
collapse 2.53\,\% $\to$ 0.94\,\%; full range adds nothing \\
\bottomrule
\end{tabular}
\caption{Adopted (d,t) arguments of \code{fitGenfitter\_a1975\_deuterium.C}.}
\label{tab:fit-dt-args}
\end{table}

Three habits of the script are worth copying. It writes to a new output directory and never overwrites a previous
production, so two productions that differ in one flag can be differenced. It checks the log of every run for the
banners of the options it set and warns if one is missing (a GENFIT without the CATIMA backend runs silently with its
own models). And it passes every argument: the header of \code{prod\_dp\_clord.sh} records that an earlier (d,p)
production passed only six arguments and silently took the macro defaults for the rest (material effects off, the H2
geometry, the baseline par, fallback on).

\section{The Unscented Kalman Filter: AtFitterUKF}
\label{sec:ukf}

\subsection{Architecture}

The UKF is built in three layers (\url{docs/development/UKF.md}):
\begin{enumerate}
\item \code{kf::UnscentedKalmanFilter} and related header-only templates in \url{AtReconstruction/AtFitter/OpenKF/}
  (vendored OpenKF, depends on Eigen3);
\item \code{kf::TrackFitterUKF} (\url{OpenKF/kalman_filter/TrackFitterUKF.h}): the physics models. The state is
  $(x, y, z, |p|, \theta, \phi)$, the measurement a 3D position, the process noise energy straggling. The process model
  propagates the state with \code{AtTools::AtPropagator} and an \code{AtRK4Stepper} through the field, with energy loss
  from an \code{AtTools::AtELossModel}; after the forward pass a Rauch--Tung--Striebel smoother runs backwards;
\item \code{EventFit::AtFitterUKF}: the \code{AtFitter} implementation.
\end{enumerate}
Unit tests are listed in \code{UKF.md}; demonstration macros are in \url{macro/tests/UKF/}.

\subsection{Construction}

\begin{verbatim}
AtFitterUKF(double charge, double mass_MeV,
            std::unique_ptr<AtTools::AtELossModel> elossModel);
\end{verbatim}
\code{charge} is in coulomb ($Z \times 1.602176634\times10^{-19}$), the mass in MeV/$c^2$. The energy-loss model is
moved into the propagator on the first fit. In every macro it is an \code{AtTools::AtELossCATIMA}; note its density
argument is in g/cm$^3$.

\subsection{What happens to one track}

\begin{enumerate}
\item Reject tracks with an invalid PRA radius; recover a NaN PRA polar angle from the track end points.
\item Reject ``beam-like'' tracks: $\theta_\mathrm{lab} = 180^\circ - $ \code{GetGeoTheta()} outside
  $[10^\circ, 170^\circ]$. The 10$^\circ$ limit (\code{fMinLabTheta}) has no setter.
\item If adaptive clustering is on (default), re-cluster the hits with a cluster radius and distance chosen from the
  $B\rho$ momentum estimate, from \code{SetClusteringParams}, or from \code{SetTargetClusters}; with
  \code{SetUseArcWalk(true)} the kNN arc walk is used instead of the $z$-seeded smoothing clusteriser.
\item Remove clusters closer than \code{SetMinClusterSpacing} (default 3\,mm) to the previous one and drop the last
  10\,\% of clusters (the Bragg end), keeping at least \code{SetMinClusters}. This step is unconditional.
\item Convert to the lab frame if \code{SetZPadPlane} $>0$.
\item Seed: position of the first cluster, momentum from $B\rho = B_z r/\sin\theta$ of the PRA circle (or
  \code{SetMomentumSeed}), direction from the cluster sequence when the lab conversion or
  \code{SetUseClusterDirSeed} is active. Initial covariance: \code{SetMeasurementSigma}$^2$ in position,
  $(\mathrm{frac}\cdot p)^2$ in momentum with \code{SetMomentumSigmaFrac}, $(1^\circ)^2$ in the angles.
\item Run the forward UKF pass and the RTS smoother, \code{SetNIterations} times; iterations after the first are
  seeded from the previous smoothed state (optionally with its covariance, \code{SetUseRefTrackWarmStart}, or with full
  relinearisation around the previous trajectory, \code{SetUseRefTrack}). A fit is abandoned after 2000\,ms, and a
  segmentation fault inside the fit is caught and the track dropped.
\item Back-extrapolate from the first smoothed state to the beam axis, either linearly along the momentum
  (default, path capped by \code{SetBackExtrapMaxPath}, default 50\,mm) or along the PRA circle
  (\code{SetUseHelixBackExtrap}). The energy lost over the path is added back as $\mathrm{d}E/\mathrm{d}x(E)\times
  \mathrm{path}$ evaluated once at the first-cluster energy.
\end{enumerate}

\subsection{Output conventions of the UKF}

\begin{itemize}
\item \code{GetKinematics()} holds the KE \emph{at the vertex} (after back-extrapolation and energy correction);
  \code{GetKinematicsXtr()} holds the KE at the first cluster. This is the reverse of \code{AtGenfitter}. The angles are
  the same in both slots and are those of the smoothed state at the first cluster, unless
  \code{SetUpdateAnglesOnBackExtrap(true)} rotates $\phi$ by the back-extrapolated arc.
\item \code{GetVertex()} is the back-extrapolated point; \code{GetTrackPropertiesStruct()} is filled (first-cluster
  position, vertex, track length from the smoothed positions, number of clusters).
\item The smoothed positions (states 1 to $N-1$) are stored in the lab frame.
\item $\chi^2$ is not the filter's $\chi^2$. It is $\sum_i |\vec r_i^{\,\mathrm{smoothed}} - \vec r_i^{\,\mathrm{meas}}|^2
  / \sigma^2$ over clusters 1 to $N-1$, with $\sigma$ from \code{SetMeasurementSigma} and
  $\mathrm{ndf}=\max(1, 3(N-1)-6)$. \code{SetChi2()} is called with $\chi^2/\mathrm{ndf}$, so for UKF output
  \code{GetChi2()} already \emph{is} the reduced value. The same is written to the event-level \code{AtFitMetadata}
  when the task provides one.
\item \code{GetGoodFit()} and \code{GetPValue()} are not set (they keep their defaults, false and 0).
\end{itemize}

\subsection{Setters}

\begin{table}[p]
\centering
\small
\begin{tabular}{@{}p{6.0cm}p{2.3cm}p{6.5cm}@{}}
\toprule
Setter & Default & Meaning \\
\midrule
\code{SetBField(XYZVector)} & $(0,0,2.85)$\,T & Signed field; see Table~\ref{tab:fit-bsign}. \\
\code{SetEField(XYZVector)} & 0 & Electric field, V/m. \\
\code{SetUKFParameters(a, b, k)} & $10^{-3}$, 2, 0 & Sigma-point scaling. \code{UKF.md}: $\alpha\ge0.1$ made sigma
points miss measurement planes. \\
\code{SetMeasurementSigma(mm)} & 1.0 & Position uncertainty of a cluster. \\
\code{SetMomentumSigmaFrac(f)} & 0.1 & Seed momentum uncertainty. \\
\code{SetMinClusters(n)} & 3 & Minimum clusters after filtering. \\
\code{SetMomentumSeed(MeV)} & off & Fixed seed momentum. \\
\code{SetNIterations(n)} & 1 & Filter passes. \\
\code{SetUseRefTrackWarmStart(b)}, \code{SetRefTrackInflation(f)} & false, 4.0 & Seed later passes with the inflated
smoothed covariance. \\
\code{SetUseRefTrack(b)} & false & Relinearise around the previous smoothed trajectory (needs $\ge2$ iterations). \\
\code{SetEnableEnergyStraggling(b)} & true & Straggling as process noise. All a1975 and 16C(p,p) macros set false. \\
\code{SetEnableMultipleScattering(b)}, \code{SetRadiationLength(mm)} & false, $7.6\times10^6$ & Highland angular
process noise. \\
\code{SetELossScaleFactor(f)} & 1.0 & Scale of the propagator \dedx{} (diagnostic). \\
\code{SetUsePerClusterCov(b)} & false & Use \code{AtHitCluster::GetCovMatrix()}. \\
\code{SetZPadPlane(mm)} & $-1$ (off) & Digi to lab conversion. \\
\code{SetAdaptiveClustering(b)} & true & Re-cluster inside the fitter. \\
\code{SetClusteringParams(r, d)} & built-in & Fixed cluster radius and distance. \\
\code{SetTargetClusters(n)}, \code{SetAdaptiveDistBounds(min, max)} & 0 (off); 4, 20 & Arc-length-aware clustering. \\
\code{SetMinClusterSpacing(mm)} & 3.0 & Minimum distance between kept clusters. \\
\code{SetUseArcWalk(b)}, \code{SetArcWalkParams(minHits, kNN)} & false; 3, 10 & kNN arc-walk clustering. \\
\code{SetUseClusterDirSeed(b)} & false & Direction seed from the clusters, sorted by distance to the axis. \\
\code{SetBackExtrapMaxPath(mm)} & 50 & Cap on the back-extrapolation; $\le0$ disables it. \\
\code{SetUseHelixBackExtrap(b)} & false & Back-extrapolate along the PRA circle. \\
\code{SetForceVertexOnBeamAxis(b)} & false & Put the vertex at $x=y=0$ (helix path). \\
\code{SetUpdateAnglesOnBackExtrap(b)} & false & Rotate $\phi$ to the vertex. \\
\code{SetBackExtrapTargetX(mm)} & NaN (off) & Straight tail to a target plane in $x$. \\
\code{SetVertexZBias(mm)} & 0 & Constant subtracted from the vertex $z$. \\
\bottomrule
\end{tabular}
\caption{\code{AtFitterUKF} setters. The options from \code{SetUseClusterDirSeed} down were added for PUMA and are
documented in \code{docs/development/PUMA\_UKF\_calibration.md}.}
\label{tab:ukf-setters}
\end{table}

\subsection{Multiple hypotheses: AtFitterUKFMulti}

\code{EventFit::AtFitterUKFMulti} holds a list of configured \code{AtFitterUKF}, one per particle:
\begin{verbatim}
void AddHypothesis(std::unique_ptr<AtFitterUKF> fitter, std::string name,
                   int chargeSign = 0);
\end{verbatim}
Every hypothesis fits its own copy of the track (the UKF re-clusters, so copies keep hypotheses independent). A
hypothesis with \code{chargeSign} $\neq0$ is skipped when the track carries a PRA charge sign that disagrees. The
fitted track with the smallest ranking value is returned and its \code{ParticleInfo.idPDG} is set to the hypothesis
name. The ranking value is \code{GetChi2()/GetNdf()} of the metadata; since \code{AtFitterUKF} already stores
$\chi^2/\mathrm{ndf}$ in the $\chi^2$ slot, the quantity compared is $\chi^2/\mathrm{ndf}^2$, which favours hypotheses
that kept more clusters. \code{SetChi2Cap(cap)} is stored but not used by \code{GetFittedTrack()} in the current code.
\code{fitUKF\_a1975.C} passes \code{chargeSign = 0} on data because the PRA charge sign was calibrated on simulation
handedness.

\subsection{UKF examples}

\paragraph{Simulation, 16C(p,p).} The macro \url{macro/Simulation/ATTPC/16C_pp/run_ukf_only.C} fits the file
\code{data/output\_digi.root} written by \code{run\_digi\_attpc.C}:
\begingroup\footnotesize
\begin{verbatim}
   double charge = 1.602176634e-19;
   double mass = 938.272;
   auto eloss = std::make_unique<AtTools::AtELossCATIMA>(3.553e-5);
   eloss->SetProjectile(1, 1, 1);
   std::vector<std::tuple<int, int, int>> mat;
   mat.push_back(std::make_tuple(1, 1, 1));
   eloss->SetMaterial(mat);

   auto ukfFitter = std::make_unique<EventFit::AtFitterUKF>(charge, mass, std::move(eloss));
   ukfFitter->SetBField(ROOT::Math::XYZVector(0, 0, 2.85));
   ukfFitter->SetUKFParameters(1e-3, 2.0, 0.0);
   ukfFitter->SetMeasurementSigma(2.0);
   ukfFitter->SetMomentumSigmaFrac(0.5);
   ukfFitter->SetEnableEnergyStraggling(false);
   ukfFitter->SetMinClusters(10);
   ukfFitter->SetNIterations(1);
   ukfFitter->SetZPadPlane(1000.0);

   AtFitterTask *fitterTask = new AtFitterTask(std::move(ukfFitter));
   fitterTask->SetInputBranch("AtPatternEvent");
   fitterTask->SetOutputBranch("AtTrackingEvent");
   fitterTask->SetFitMetadataBranch("AtFitMetadata");
   fitterTask->SetPersistence(kTRUE);
\end{verbatim}
\endgroup
\code{UKF.md} records, for digitised 16C(p,p) at $20^\circ$ c.m., 98\,\% convergence, a momentum bias of $-1.6$\,\%
and 1.0\,\% RMS; part of the bias is a 5\,\% difference between CATIMA and Geant4 stopping powers at 0.5\,MeV, which
\code{SetELossScaleFactor(1.08)} removes on simulation.

\paragraph{Data, a1975.} \url{macro/Unpack_HDF5/a1975/UKF/pipeline/fitUKF_a1975.C} builds one hypothesis per name in a
comma-separated list (proton, deuteron, triton, He3, alpha) and combines them in \code{AtFitterUKFMulti}:
\begingroup\footnotesize
\begin{verbatim}
   auto eloss = std::make_unique<AtTools::AtELossCATIMA>(gasDensity);
   eloss->SetProjectile(p.A, p.Z, massAmu);
   ...
   auto ukf = std::make_unique<EventFit::AtFitterUKF>(charge, p.massMeV, std::move(eloss));
   ukf->SetBField(ROOT::Math::XYZVector(0, 0, bFieldSign * bFieldMag));
   ukf->SetUKFParameters(1e-3, 2.0, 0.0);
   ukf->SetMeasurementSigma(measSigma);
   ukf->SetMomentumSigmaFrac(momSigmaFrac);
   ukf->SetEnableEnergyStraggling(false);
   ukf->SetMinClusters(minClusters);
   ukf->SetNIterations(nIter);
   ukf->SetUseRefTrack(refTrack);
   ukf->SetRefTrackInflation(refInflation);
   ukf->SetZPadPlane(1000.0);
\end{verbatim}
\endgroup
with defaults \code{bFieldSign = -1}, \code{bFieldMag = 2.85}, \code{gasDensity = 9.0e-5}, \code{measSigma = 2.0},
\code{momSigmaFrac = 0.3}, \code{nIter = 1}, \code{minClusters = 10}. It reads \code{<run>\_reco.root} and writes
\code{<run>\_ukf<suffix>.root}:
\begin{verbatim}
root -b -q 'pipeline/fitUKF_a1975.C("run_0116", -1, "proton", -1)'
\end{verbatim}
\code{GENFIT\_PROGRESS.md} records that multiple scattering changes the median KE by about 1\,\% in H2 at 1\,bar and
that lowering \code{minClusters} from 10 to 4 recovers 29\,\% more fitted tracks on run\_0106.

\paragraph{PUMA.} \code{docs/development/PUMA\_UKF\_calibration.md} documents the opt-in options of
Table~\ref{tab:ukf-setters} (cluster-direction seed, arc walk, helix back-extrapolation, target cluster count,
\code{SetVertexZBias(8.6)} for a constant PSA $z$ offset) and the configuration block used in
\url{macro/Simulation/ATTPC/PUMA/run_digi_attpc.C}.

\paragraph{Display.} \code{AtUKFDisplay} (\url{AtReconstruction/AtFitter/AtUKFDisplay.h}) is a stand-alone TCanvas
display with live refitting that does not need OpenGL.

\section{Other fitters}
\label{sec:fit-other}

\paragraph{Legacy AtGenfit.} \code{AtFITTER::AtGenfit} derives from the deprecated \code{AtFitterOld} and runs in the
deprecated \code{AtFitterTaskOld}; all three are compiled only with GENFIT2, and
\url{AtReconstruction/CMakeLists.txt} lists the last two as ``deprecated, to be removed''. Its constructor is
\code{AtGenfit(magfield, minbrho, maxbrho, eLossFile, gasMediumDensity, pdg = 2212, minit = 5, maxit = 20,
noMatEffects = kFALSE)}; it carries its own merging, single-vertex and reclustering options
(\code{EnableMerging}, \code{EnableSingleVertexTrack}, \code{EnableReclustering(on, radius, size)}) and a
\code{SetSimulationConvention(bool)} flag in place of a signed field. The a1975 macro
\url{macro/Unpack_HDF5/a1975/UKF/pipeline/fitGenfit_a1975.C} shows its use. Use it only to reproduce older results.

Both GENFIT fitters build their measurements with \code{genfit::AtSpacepointMeasurement}, which now uses the covariance
stored in the cluster when one is set (as \code{AtGenfitter} does) and otherwise a fixed $(2\,\mathrm{mm})^2$,
$(2\,\mathrm{mm})^2$, $0.08$\,cm$^2$ diagonal. The comment there records that an earlier 10\,mm placeholder made the fit
follow the seed.

\paragraph{Monte Carlo fitters.} \code{MCFitter::AtMCFitter} (\url{AtReconstruction/AtFitter/AtMCFitter.h}) samples
parameters from \code{AtParameterDistribution}s, simulates and digitises each trial event with \code{AtClusterize},
\code{AtPulse} and an optional PSA, and keeps the trials with the lowest objective function; derived classes implement
\code{CreateParamDistros()}, \code{SetParamDistributions()}, \code{ObjectiveFunction()} and \code{SimulateEvent()}.
\code{AtMCFission} is the implementation for fission events. \code{AtMCFitterTask} drives them from
\code{AtPatternEvent} (\code{SetPatternBranchName}). These classes are not used in the a1975/a1954 analyses and are not
described further here.

\paragraph{AtMCQMinimization.} \code{MCMinimization::AtMCQMinimization}
(\url{AtReconstruction/AtMinimization/AtMCQMinimization.h}) fits a single track by simulating it (Euler integration
of the Lorentz force with energy loss, electrons spread over the pad plane with the charge response of
\code{AtDigiPar}) and comparing the pad charge pattern with the data in a random walk with shrinking steps. The
objective combines a charge term with either a position term (default) or a range term (\code{SetUseRangeChi2}). It is
not an \code{AtFitter} and has no task; it is called from a macro with a \code{TrackSeed} and the hits:
\begin{verbatim}
   auto eLoss = std::make_shared<AtTools::AtELossCATIMA>(density, material);
   eLoss->SetProjectile(A, Z, massAmu);
   auto min = MCMinimization::AtMCQMinimization();
   min.SetMap(fMap);          // pad plane geometry, mandatory
   min.SetParticle(A, Z);
   min.SetELossModel(eLoss);
   min.SetBFieldSign(-1);     // the par file carries no sign
   min.Minimize(seed, hits);
   auto res = min.GetFitPar();
\end{verbatim}
\url{macro/tests/AT-TPC/MCQMinimizationOnData.C} applies it to a1975 run\_0106 and reports, over the first 35 events,
a mean difference to \code{AtGenfitter} of $-1.0^\circ$ in $\theta$ and $+0.57$\,MeV in energy; with the wrong field
sign the mean hit-to-pad distance goes from about 2\,mm to about 35\,mm.

\section{Comparing fitters and tuning}
\label{sec:fit-compare}

All fitters read the same \code{AtPatternEvent}, so a difference between two fits of one reco file comes from the
fitter alone. Table~\ref{tab:fit-compare} summarises the differences between the two maintained fitters and the legacy
one.

\begin{table}[h]
\centering
\small
\begin{tabular}{@{}p{2.4cm}p{3.9cm}p{4.0cm}p{4.5cm}@{}}
\toprule
 & \code{AtFitterUKF} & \code{AtGenfitter} & \code{AtGenfit} (legacy) \\
\midrule
Algorithm & Unscented KF, RTS smoother & GENFIT reference-track KF & GENFIT reference-track KF \\
Energy loss & \code{AtELossModel} (CATIMA) & none, Bethe--Bloch, table or CATIMA (Section~\ref{sec:gf-eloss}) & table file \\
Iterations & \code{SetNIterations} (1) & min--max (2--5) & min--max (5--20) \\
Field sign & signed \code{SetBField} & signed field & \code{SetSimulationConvention} \\
Re-clustering & inside the fit (default on) & none; opt-in merge & merge, recluster options \\
Species selection & \code{AtFitterUKFMulti} & \code{SetPIDGate} & none \\
Task & \code{AtFitterTask} & \code{AtFitterTask} & \code{AtFitterTaskOld} \\
\bottomrule
\end{tabular}
\caption{The three Kalman-filter fitters side by side.}
\label{tab:fit-compare}
\end{table}

\paragraph{Both fitters in one job.} \url{macro/Unpack_HDF5/a1975/UKF/pipeline/fitBoth_a1975.C} runs \code{AtGenfitter}
and \code{AtFitterUKF} with the proton hypothesis on the same events and writes two tracking branches,
\code{AtTrackingEventGenfit} and \code{AtTrackingEventUKF}, which \code{dispBoth\_a1975.C} and
\code{run\_eve\_both\_a1975.C} overlay on the hits:
\begingroup\footnotesize
\begin{verbatim}
root -b -q 'pipeline/fitBoth_a1975.C("run_0106", 200, "/mnt/f/a1975/reco/", "/tmp/")'
\end{verbatim}
\endgroup

\paragraph{Tuning macros.} The directory \url{macro/Unpack_HDF5/a1975/UKF/pipeline/} holds the drivers of the studies
quoted in this chapter: \code{featureScan\_a1975.C} (\code{SetMeasSigma}, \code{SetDiffusion} against median
$\chi^2/\mathrm{ndf}$ and good-fit rate), \code{testMerge\_a1975.C}, \code{mergeQuality\_a1975.C} and
\code{mergeDiag\_a1975.C} (continuity merging), \code{ukfMSscan\_a1975.C} and \code{ukfMinClus\_a1975.C} (UKF multiple
scattering and minimum clusters). \url{macro/Unpack_HDF5/a1975/UKF/README.md} records that \code{measSigma = 0.5}\,mm with
\code{momSigmaFrac = 0.1} gives about 26\,\% better elastic-peak FWHM than the 2.0 and 0.3 defaults of
\code{fitUKF\_a1975.C}. Such knobs depend on the channel and the field: re-validate them for each reaction, and write
each variant to its own output file (\code{outSuffix}).

\section{The fit output}
\label{sec:fit-output}

\subsection{Classes}

\code{AtTrackingEvent} (one per event, in a \code{TClonesArray} of size one) holds the event ID, the copies of the input
\code{AtTrack}s (\code{GetTrackArray()}, returned by value) and the fitted tracks (\code{GetFittedTracks()}, a vector
of \code{unique\_ptr<AtFittedTrack>}). A fitted track keeps the ID of the pattern track it came from
(\code{GetTrackID()}); with the \code{AtTrackFinderTC} finder that ID is a cluster label, not an array index, so
fitted tracks are matched to pattern tracks and to \code{AtPIDEvent} entries by ID.

\code{AtFittedTrack} (Table~\ref{tab:fit-fittedtrack}) stores per-particle vectors; both fitters fill index 0.

\begin{table}[h]
\centering
\small
\begin{tabular}{@{}p{5.2cm}p{4.8cm}p{4.8cm}@{}}
\toprule
Accessor & \code{AtGenfitter} & \code{AtFitterUKF} \\
\midrule
\code{GetKinematics()} (KE MeV, $\theta$, $\phi$ rad) & raw fit at the first measurement point & at the vertex, after
back-extrapolation \\
\code{GetKinematicsXtr()} & after \code{SetBackExtrapToAxis}/\allowbreak\code{SetManualELoss}; equals the raw fit if neither is
on & at the first cluster \\
\code{GetVertex()} (mm, lab frame) & final state position (axis point or first fitted point) & back-extrapolated point \\
\code{GetSmoothedPositions()} & fitted position at each measurement, lab frame & smoothed states 1..$N-1$, lab
frame \\
\code{GetParticleInfo()} & \code{idPDG} = $Z$ as a string, charge $Z$, mass (u) & same; \code{AtFitterUKFMulti}
writes the hypothesis name into \code{idPDG} \\
\code{GetTrackPropertiesStruct()} & not filled (defaults $-1$) & filled \\
\code{GetTrackMetadata()} & see below & see below \\
\bottomrule
\end{tabular}
\caption{What each fitter writes into \code{AtFittedTrack}.}
\label{tab:fit-fittedtrack}
\end{table}

\code{AtFitTrackMetadata} (\code{GetTrackMetadata()}) carries:
\begin{itemize}
\item \code{GetChi2()}, \code{GetNdf()}: for \code{AtGenfitter} GENFIT's total $\chi^2$ and ndf; for the UKF the
  reduced $\chi^2$ and its ndf. The same macro cannot read both as the same quantity.
\item \code{GetFitConverged()}; \code{GetGoodFit()} (\code{AtGenfitter} quality flag only);
  \code{GetPValue()} (always 0);
\item \code{GetMatEffects()}, \code{GetMatEffectsFallback()} (\code{AtGenfitter} only): whether material effects were
  active in the fit that produced the result, and whether it came from the no-material retry.
\end{itemize}
The event-level \code{AtFitMetadata} branch is filled by the UKF; \code{AtGenfitter} ignores it, so in GENFIT output the
branch exists but holds no track entries.

\subsection{Frame of the smoothed positions}

Both fitters return smoothed positions in the lab frame, $z_\mathrm{lab} = z_\mathrm{pad} - z_\mathrm{digi}$. The hits
and clusters stored in the same file (in \code{AtPatternEvent} and in \code{AtTrackingEvent::GetTrackArray()}) are in
the digi frame. To overlay a fit on its clusters, transform one of them: with $z_\mathrm{pad}=1000$\,mm a smoothed point
and the cluster it belongs to satisfy $z_\mathrm{fit} + z_\mathrm{cluster} \approx 1000$\,mm, and $x$, $y$ agree
directly.

\subsection{Reading the output}

The recommended loop starts from the PID results, skips entries without them, and matches fitted tracks by ID. This
is the core of \url{macro/Unpack_HDF5/a1975/D2_UKF/ab_catima_collapse.C}:
\begin{verbatim}
   TTree *t = (TTree *)fl->Get("cbmsim");
   TClonesArray *te = nullptr, *pe = nullptr;
   t->SetBranchAddress("AtTrackingEvent", &te);
   t->SetBranchAddress("AtPIDEvent", &pe);
   for (Long64_t i = 0; i < t->GetEntries(); ++i) {
      t->GetEntry(i);
      if (!te || !pe || te->GetEntries() == 0 || pe->GetEntries() == 0)
         continue;
      auto *ev = (AtTrackingEvent *)te->At(0);
      auto *pidev = (AtPIDEvent *)pe->At(0);
      if (!ev || !pidev)
         continue;
      std::map<int, AtFittedTrack *> fmap;
      for (auto &ft : ev->GetFittedTracks())
         if (ft)
            fmap[ft->GetTrackID()] = ft.get();
      for (auto &sr : pidev->GetSpyral()) {
         if (!sr.valid || sr.brho <= 0)
            continue;
         auto it = fmap.find(sr.trackID);
         if (it == fmap.end())
            continue;
         const auto &m = it->second->GetTrackMetadata();
         const bool isBad = !(m->GetNdf() > 0 && m->GetChi2() > 0);
         ...
\end{verbatim}
Analysis macros then read \code{GetKinematicsXtr()} for GENFIT output (for example
\code{ex\_dt\_a1975.C}, ``back-extrapolated: the value the analysis uses'').

\section{Pitfalls}
\label{sec:fit-pitfalls}

\begin{enumerate}
\item \textbf{Repeated entries in the output tree.} When the pattern recognition cannot handle an event it leaves
  \code{AtPatternEvent} empty (\code{AtPRAtask} logs ``pattern recognition returned nothing for this event''). For such
  an event \code{AtFitterTask::Exec()} returns before clearing its arrays, so the tree entry repeats the previous
  event's \code{AtTrackingEvent}, with the previous event ID. \code{AtPIDTask} clears first, so the same entry has an
  empty \code{AtPIDEvent}. Loops over \code{GetFittedTracks()} alone count these tracks twice. Loop over the PID (or
  Spyral) results and match fitted tracks by track ID, as in the example above, or reject entries whose event ID
  repeats.
\item \textbf{Zero energy loss below $\beta\gamma=0.05$.} With material effects on and no \dedx{} source for that
  region, GENFIT applies no stopping power below KE 1.17\,MeV (proton) or 3.5\,MeV (triton). Load a hybrid table and/or
  enable \code{SetCatimaELoss}.
\item \textbf{Parameterisation-only table.} Passing a table without \code{SetELossHybrid(kTRUE, rho)} puts GENFIT in
  table-only mode, which throws above the table's highest energy.
\item \textbf{Collapsed fits are kept.} A GENFIT fit can return with $\mathrm{ndf}\le0$ and still carry kinematics.
  Require $\mathrm{ndf}>0$ (or \code{GetGoodFit()}) before using a track.
\item \textbf{Navigator state.} Material-effects productions made before the per-track navigator re-seed collapsed in
  contiguous blocks; they are not comparable with later productions.
\item \textbf{Vertex gap counted twice.} \code{SetManualELoss} together with material effects and
  \code{SetBackExtrapToAxis} adds the gap loss twice ($+5.25$\,\% on the (d,t) low branch).
\item \textbf{Kinematics slots.} \code{GetKinematics()} is the raw first-point fit for \code{AtGenfitter} and the
  vertex value for the UKF; \code{GetKinematicsXtr()} is the reverse. Check which fitter produced a file before
  reading it.
\item \textbf{$\chi^2$ conventions.} \code{GetChi2()} is total for GENFIT and reduced for the UKF. At the GENFIT default
  $\sigma_0=4$\,mm the residuals are about 0.6\,mm, $\chi^2/\mathrm{ndf}$ is of order 0.1 and grows with track length,
  so a fixed cut selects long tracks out. Set any threshold from the measured distribution.
\item \textbf{Mirrored frame of smoothed positions.} Smoothed positions are in the lab frame, clusters in the digi
  frame: $z_\mathrm{fit} + z_\mathrm{cluster} = z_\mathrm{pad}$.
\item \textbf{Field sign.} The sign is a property of the setup; forward tracks hide a wrong sign. See
  Table~\ref{tab:fit-bsign}.
\item \textbf{Simulated polar angles.} With a fit sign that matches the generated field, the fitted polar angle of a
  simulated track compares directly with the true one. Measure this for each simulation.
\item \textbf{z sort near $90^\circ$.} Sorting clusters by $z$ scrambles tracks close to $90^\circ$;
  \code{SetUseClusterOrder} is the alternative but lost fits in the tests made so far and is not adopted.
\item \textbf{Fallback mixes models.} With \code{SetMatEffectsFallback(kTRUE)} (the default) a material-effects
  production contains no-material fits; filter on \code{GetMatEffectsFallback()} or turn the fallback off.
\item \textbf{Geometry is the material.} With material effects on, the medium of the loaded geometry must be the target
  gas; \code{ATTPC\_H1bar\_geomanager.root} is ordinary hydrogen.
\item \textbf{Silent defaults.} The fit macros have long positional argument lists; a call that stops early takes the
  defaults for everything after it. Pass every argument in production scripts and check the log banners.
\item \textbf{PID gate drawn on another plane.} The internal gate uses \code{AtSpyralPID} with 30 minimum points, or
  \code{AtHelix3D} with \code{SetPIDGateFrom3D}, with the gains of \code{SetDedxGain}; a gate drawn under other settings
  selects a different sample. A missing gate file disables gating with only a log message.
\end{enumerate}
